\section{Introduction}
Algebraic simplification involves using algebraic laws to normalise
expressions with unknowns. For example, the commutative monoid
axioms---associativity, neutrality, and commutativity---over
integers with addition serve to
simplify the left hand expression to the right hand expression below:
%
\[-6 + (x+3)+(y+x)
\;{}\xmapsto{\text{simplify}}{}\;
-3+2x+y
\]
%
%% Despite the suggestive name, simplification does not always make
%% terms simpler; rather, it chooses a canonical representative.
Many application domains make use of this kind of simplification.
%
For example, algebraic simplification is often a useful first step in
a program optimiser, to avoid the need to analyse and transform
distinct but equivalent programs.
%
The present work focuses on another application domain: theorem
provers and programming languages based on type theory.
%
In these systems, users often need to prove to a type checker that
terms are equivalent, but constructing the proofs often involves rote
algebraic simplification steps that users resent having to produce
manually.

To free users from the need to construct rote algebraic proofs,
de\-pen\-dent\-ly ty\-ped languages and their ecosystems often
include simplifiers for common algebraic structures such as
monoids, semi-rings, and rings.
%
With these simplifiers, users need only establish the structures'
axioms, such as neutrality, associativity and commutativity, and can
then call the simplifiers to discharge rote simplification steps.
%
Implementation strategies for simplifiers vary:  some,
such as the Rocq ring solver~\cite{coq-ref-man-8:ring-solver},
use tactics to simplify  algebraic terms in typing goals,
while others, such as the Agda ring solver~\cite{kidney:ring-solver},
use proof-by-reflection to construct propositions to discharge
equations.
Some simplifiers, such as the Rocq ring solver, group together several algebraic
structures, while others, such as the system by \citet{gregoire-mahboubi:commutative-ring-done-right}, generalise several distinct structures to one
structure.
The
state-of-the-art are standalone simplifiers that, through heuristics
and long-term development, can deal with common cases.

\subsection{Representation Theorems}

Universal algebra has a long tradition concerning algebraic
simplification under the collective name `representation
theorems'. Each such representation theorem characterises canonical
representatives of algebraic expressions in terms of (typically inductive)
constructions such as reduced-words and formal polynomials. Such
characterisations often
reuse existing representation theorems of simpler algebraic
structures or familiar algebraic structures such as the integers or
the natural numbers.

The present work makes use of two kinds of representation theorems.
%
For a free algebra (abbreviated \emph{fral}), a representation theorem
amounts to an algebraic structure that interprets the algebraic
operations to produce a single canonical representative for all input
terms that are equivalent according to the algebraic laws.
%
For a free extension (abbreviated \emph{frex}), a representation
theorem chooses a canonical representative using the algebraic laws
while also evaluating concrete elements.
%
Free algebras and free extensions are related: each free extension is
also a free algebra of a theory specialised to a concrete algebra by
adding axioms over the concrete elements.

For example, for commutative monoids the fral representation theorem
states that the free commutative monoid over $n$ variables is
represented by the set of $n$-tuples of naturals $\naturals^n$.
We simplify with this representation by evaluating
a term in the fral and then reifying it back as a term:
\[
-6 + (x+3) + (y+x)
{}\xmapsto{\text{evaluate}(x_0 + (x_2 + x_1) + (x_3 + x_2))}{}
(1, 1, 2, 1)
{}\xmapsto{\text{reify}(x_0 \mapsto -6, x_1 \mapsto 3, x_2 \mapsto x, x_3 \mapsto y)}{}
-6 +3 + 2x + y
\]\label{new-intro fral example}
%
In the fral there is no concept of concrete elements, so fral
simplification must treat constants such as $-6$ and $3$ in the same
way as variables such as $x$ and $y$, as abstract and distinct
indeterminates.

The frex representation theorem for commutative monoids
states that the free extension of a commutative monoid over $C$ by $n$
variables is represented by the set $C \times \naturals^n$, where
the element $(c, a_1,
\ldots, a_n)$ represents the expression $c + a_1x_1 + \ldots +
a_nx_n$.
%
Simplifying a term using the frex representation involves evaluating a
term in the frex and then reifying it back as a term:
\[
-6 + (x+3) + (y+x)
{}\xmapsto{\text{evaluate}_{\integers}}{}
(-3, 2, 1)
{}\xmapsto{\text{reify}}{}
-3 + 2x + y
\]
The frex representation distinguishes variables from concrete
elements, gathering the latter together and evaluating them using the
operations of the concrete commutative monoid in use.

Both kinds of representation theorem allow us to choose a unique
syntactic representative, i.e., a normal form.
Such normalization by evaluation and then reification
also applies to more sophisticated notions of algebra that include the
equational theories of $\lambda$-calculi, and is familiar in those
settings as \emph{normalization-by-evaluation}~\cite{berger:nbe}.  Its first systematic
applications were in the formal study of various simply typed
calculi~\cite{altenkirch-et-al:cat-reconstruction,%
  cubric-dybjer-scott:normalization-and-the-yoneda-embedding,%
  altenkirch-et-al:nbe-stlc-sums}
and category theoretic
constructions~\cite{beylin-dybjer:extracting-coherence-for-monoidal-cats}.
It has served as a conversion-checking technique during the
type-checking of dependently typed calculi from their
inception~\cite{martin-lof:an-intuitionistic-theory-of-types},
gaining adoption after the seminal works
of~\citet{abel-aehlig-dybjer:nbe-ml-tt-universe,abel-coquand-dybjer:nbe-mltt-eq},
even for sophisticated
calculi~\cite{DBLP:journals/pacmpl/0001VW17,sterling-angiuli:normalization-for-cubical-type-theory,DBLP:journals/jfp/HuJP23}.

This article describes an extensible dependently typed library for
algebraic simplifiers based on fral and frex representation theorems,
building on our previous work on free extensions for
partial evaluation~\cite{yallop-et-al:partially-static-data-as-frex}.
%
Current implementations of algebraic simplifiers,
even in dependently typed settings, are restricted to
implementing the computational representation---i.e.~the data-structures
needed for the normal form together with the
normalisation-by-evaluation algorithm---alongside a formalisation of
the soundness proofs.
%
This work investigates what can be gained by encoding, in addition,
the full meta-theory of these representation
theorems, including generic representations of theories, their
algebras and algebra homomorphisms, and the universal properties of
the fral and frex.
%
For simplicity of development and exposition we apply this generic
machinery to a handful of familiar monoid varieties.
%
Moreover, we ensure all of these concepts remain computational by avoiding
the temptation of postulating axioms that could hinder reduction of closed
terms.
This last task is challenging to combine with
interactive performance, as formalising universal properties tends to
produce large terms that slow
type-checkers~\cite{gross-chlipala-spivak:performant-cat-theory-in-coq}.


\subsection{Paper Outline and Contributions}

\secref{new-overview} present background material: a
review of the mathematical foundation for \Frexlib{},
and a brief \idris{} tutorial that reviews
setoid-based equational reasoning.

\secsref{frexlib}{core} present our central contribution, a
fundamentally new approach to building algebraic simplifiers.
%
The standard existing approach is to write an ad-hoc simplifier for
some particular algebraic structure such as rings.
We take a different approach, teaching the implementation the basic concepts of
universal algebra --- signatures, theories and models, homomorphisms,
and universal properties (\secref{frexlib}) --- then build a
completely generic solver based on free algebras and free extensions
that can be instantiated with a particular algebra to discharge
concrete proof obligations (\secref{core}).  This new approach is
inherently modular and extensible, and delivers solvers that are sound
and complete by construction.  We have created implementations of our design in
two dependently typed languages: \Fragmentlib{} in \agda{},
and \Frexlib{}
in \idris{}\footnote{
{Available here:
  \url{https://github.com/frex-project/agda-fragment}
and here:
  \url{https://github.com/frex-project/idris-frex}}}.

\secref{certification} explains the completeness guarantees of the
library, and covers proof extraction, simplification, pretty-printing
and certification. The technical \secsref{core}{certification}
are aimed at library designers, and may be skimmed or skipped at
first reading.

\secref{reflection} considers a natural question: can one use
reflection to invoke \Frexlib{} automatically? The answer is a
qualified `yes', requiring much library-developer effort, but leading
to real advantages in \agda{} and limited advantages in \idris{}.

\secref{evaluation} reports some supplementary evaluation of
\Frexlib{}.  The key properties of our design are guaranteed by the
type theories of the languages in which we realise it: it delivers
sound and complete solvers in a completely generic way, with support
for proof extraction, certification, etc.  However, the practical
questions of usability and viability for interactive development
cannot be established by theorems and so we have also carried out some
experiments.  These experiments focus on the varieties of monoids that
also serve as our running example, and establish that the generic
solver is comfortably fast enough for interactive use, and can be
extended with new algebras in a modular way and without enormous
effort.

\secref{lessons} discusses system design issues that we encountered
with \Frexlib{}, and \secsref{related}{conclusion} conclude with
related and further work.

Appendices~\ref{app:first-appendix}--\ref{app:last-appendix}
have more information about \Frexlib's
\FrexlibLoC{}-line \idris{} codebase, and example code extraction.
%%
%% We include only code excerpts that convey the
%% key ideas; we refer to the implementation for full
%% details.

% Short contributions.
% Structure of paper
